\chapter{Podsumowanie i wnioski}
\label{cha:summary}

Celem niniejszej pracy magisterskiej było zaprojektowanie i zaimplementowanie nowoczesnego narzędzia do generowania syntetycznych danych relacyjnych, wspieranego przez modele sztucznej inteligencji. Cel ten został osiągnięty poprzez stworzenie kompletnej aplikacji webowej, która integruje szybkość tradycyjnych generatorów z kreatywnością modeli \gls{LLM}.

\section{Osiągnięcia}
W ramach prac zrealizowano wszystkie założone wymagania funkcjonalne i niefunkcjonalne:
\begin{itemize}
    \item \textbf{Hybrydowy Silnik Generujący}: Stworzono elastyczny mechanizm łączący bibliotekę Faker, szablony Jinja2 oraz modele \gls{LLM} (OpenAI/Ollama). Pozwala to na generowanie danych o dowolnym stopniu skomplikowania – od prostych typów liczbowych po złożone opisy tekstowe.
    \item \textbf{Obsługa Relacji}: Zaimplementowano algorytm automatycznego rozwiązywania zależności między tabelami, co gwarantuje spójność referencyjną generowanych zbiorów danych.
    \item \textbf{Wydajność}: Dzięki zastosowaniu architektury asynchronicznej (Celery + Redis), system jest w stanie obsłużyć generowanie tysięcy rekordów bez blokowania interfejsu użytkownika.
    \item \textbf{Konteneryzacja}: Całość rozwiązania została zamknięta w kontenerach Docker, co drastycznie upraszcza proces wdrożenia i konfiguracji w dowolnym środowisku.
\end{itemize}

\section{Dalszy rozwój aplikacji}
Zrealizowany projekt stanowi solidną podstawę do dalszego rozwoju. Zidentyfikowano kilka obszarów, w których system mógłby zostać rozbudowany:

\begin{enumerate}
    \item \textbf{Wsparcie dla kolejnych silników bazodanowych}: Obecnie system wspiera eksport do PostgreSQL. Rozszerzenie o natywne wsparcie dla MySQL, Oracle oraz baz NoSQL (np. MongoDB) zwiększyłoby uniwersalność narzędzia.
    \item \textbf{Generowanie na podstawie istniejącej bazy}: Ciekawą funkcjonalnością byłaby inżynieria wsteczna - podłączenie się do istniejącej bazy danych, automatyczne odczytanie jej schematu i wygenerowanie danych testowych ,,podobnych'' do danych produkcyjnych (z zachowaniem anonimizacji).
    \item \textbf{Inteligentne wykrywanie typów}: Wykorzystanie \gls{LLM} do automatycznej sugestii typów danych na podstawie samej nazwy kolumny (np. wpisując nazwę "pesel", system sam proponuje odpowiedni generator regex).
\end{enumerate}

Podsumowując, stworzone narzędzie wypełnia lukę na rynku generatorów danych testowych, oferując unikalną możliwość generowania treści semantycznie poprawnych i kontekstowych, co jest kluczowe w dobie nowoczesnych aplikacji opartych na danych.
